% Ampliación sustantiva de la completitud coinductiva Hodge.
% Módulo destinado al capítulo Hodge de la Parte IV.

\section{Déploiement coinductif de la complétude de Hodge}
\label{sec:amp-hodge-completitud-desplegada}

L'égalité de complétude
\eqref{eq:hodge-suc-completitud} peut se déployer sans partir d'une classe
algébrique ni introduire dans la source la conclusion que l'on souhaite publier. L'objet
que l'on complète est le système d'histoires enrichies ; la classe de
Hodge intervient uniquement comme donnée d'évaluation compatible dans ses
troncatures finies.

\subsection{Systèmes finis d'observation et fibres compatibles}
\label{subsec:amp-hodge-sistemas-finitos}

Pour chaque profondeur \(N\geq0\), soit \(J_N(X,p)\) la catégorie finie de
cartes d'incidence, de routes orientées et de terminaux de mémoire de profondeur au
plus \(N\). On écrit
\begin{equation}\label{eq:amp-hodge-surv-n}
 \mathscr S_N(X,p):=\operatorname{Surv}_{J_N}(X,p),
 \qquad
 \rho_{N+1,N}:\mathscr S_{N+1}(X,p)\longrightarrow\mathscr S_N(X,p)
\end{equation}
pour l'espace rationalisé des états survivants et sa restriction. Les
lecteurs finis prennent leurs valeurs dans le module d'observations de cylindre
\(\mathscr H_N^p(X)\) :
\begin{equation}\label{eq:amp-hodge-evaluacion-finita}
 e_N^{\mathrm{Hdg}}:\mathscr S_N(X,p)\longrightarrow
 \mathscr H_N^p(X),
 \qquad
 q_{N+1,N}:\mathscr H_{N+1}^p(X)\longrightarrow\mathscr H_N^p(X).
\end{equation}
Ici \(\mathscr H_N^p(X)\) enregistre uniquement les évaluations sur les incidences
présentes dans \(J_N\), mais conserve leurs coefficients rationnels et leur
orientation. Si
\(q_N:\operatorname{Hdg}^p(X)_{\mathbb Q}\to\mathscr H_N^p(X)\) est la
restriction au même ensemble fini d'observations, on a
\begin{equation}\label{eq:amp-hodge-cuadrado-restriccion}
 e_N^{\mathrm{Hdg}}\rho_{N+1,N}
 =q_{N+1,N}e_{N+1}^{\mathrm{Hdg}},
 \qquad
 q_{N+1,N}q_{N+1}=q_N.
\end{equation}

Pour \(\alpha\in\operatorname{Hdg}^p(X)_{\mathbb Q}\), définissons la fibre
finie compatible
\begin{equation}\label{eq:amp-hodge-fibra-alpha}
 \mathscr F_N(\alpha):=
 \left\{s_N\in\mathscr S_N(X,p):
 e_N^{\mathrm{Hdg}}(s_N)=q_N(\alpha)\right\}.
\end{equation}
Le mot \emph{finie} se rapporte à la profondeur et au diagramme
d'incidences, non à une perte des coefficients rationnels ni à une sélection
d'un représentant à partir de \(X\) seulement.

\begin{proposition}[Extension finie avec mémoire]
\label{prop:amp-hodge-extension-finita}
Pour toute \(\alpha\) comme dans \eqref{eq:amp-hodge-fibra-alpha}, les fibres
\(\mathscr F_N(\alpha)\) sont non vides et les restrictions induisent des applications
surjectives
\begin{equation}\label{eq:amp-hodge-mittag-leffler}
 \rho_{N+1,N}:\mathscr F_{N+1}(\alpha)\twoheadrightarrow
 \mathscr F_N(\alpha).
\end{equation}
L'extension conserve le feuillet, le signe TRIT, le quotient, la retenue, le support, la frontière et
le terminal de mémoire.
\end{proposition}

\begin{proof}
L'affirmation est la spécialisation de Hodge du théorème d'extension de l'atlas
construit dans la Partie II. Il convient d'en rendre le mécanisme visible. Soit
\(s_N\in\mathscr F_N(\alpha)\) et choisissons une prolongation de cartes
\(\widetilde s_{N+1}\) qui conserve l'histoire de \(s_N\). Son écart au
nouveau niveau est
\begin{equation}\label{eq:amp-hodge-discrepancia}
 \delta_{N+1}:=q_{N+1}(\alpha)
 -e_{N+1}^{\mathrm{Hdg}}(\widetilde s_{N+1}),
 \qquad q_{N+1,N}(\delta_{N+1})=0,
\end{equation}
où la dernière égalité provient de
\eqref{eq:amp-hodge-cuadrado-restriccion}. Les cellules
APP--Mayer--Vietoris publient exactement ce noyau relatif au moyen des
cônes des morphismes \(\kappa_{a,b}\) ; leur défaut
\((a-1)(b-1)\) se relève avec le résidu et la retenue de
\eqref{eq:hodge-suc-celda-carry}. D'après le théorème d'extension finie, il existe une
correction relative \(\eta_{N+1}(\delta_{N+1})\), à support dans les cartes
nouvelles, telle que
\begin{equation}\label{eq:amp-hodge-extension-operativa}
 \rho_{N+1,N}\eta_{N+1}(\delta_{N+1})=0,
 \qquad
 e_{N+1}^{\mathrm{Hdg}}\eta_{N+1}(\delta_{N+1})=\delta_{N+1}.
\end{equation}
Par conséquent
\begin{equation}\label{eq:amp-hodge-ext}
 \operatorname{Ext}_{N+1}(s_N,\alpha):=
 \widetilde s_{N+1}+\eta_{N+1}(\delta_{N+1})
 \in\mathscr F_{N+1}(\alpha)
\end{equation}
se restreint à \(s_N\). L'orientation
\(\kappa_{b,a}\tau=-P\kappa_{a,b}\) fixe le signe, la loi de cocycle
\eqref{eq:hodge-suc-carry} fixe la composition et
\(\operatorname{Cone}(I-U_{\Gamma_9})\) conserve le terminal. Le niveau initial
s'obtient à partir de l'atlas basé de germes. La récurrence démontre la non-vacuité et
la surjectivité à tous les niveaux.
\end{proof}

\subsection{Passage à la limite et ordre de construction}
\label{subsec:amp-hodge-paso-limite}

\begin{theorem}[Complétude coinductive déployée]
\label{thm:amp-hodge-completitud-coinductiva}
Avec les données déjà construites dans la Partie II --- les systèmes
\eqref{eq:amp-hodge-surv-n}--\eqref{eq:amp-hodge-evaluacion-finita}, la
compatibilité \eqref{eq:amp-hodge-cuadrado-restriccion}, l'extension finie
de la proposition \ref{prop:amp-hodge-extension-finita}, la séparation des
évaluations de cylindre et la conservation du terminal ---, pour chaque
\(\alpha\in\operatorname{Hdg}^p(X)_{\mathbb Q}\), il existe une histoire
enrichie
\begin{equation}\label{eq:amp-hodge-xi-limite}
 \xi_\alpha=(s_0,s_1,s_2,\ldots)
 \in\varprojlim_N\mathscr F_N(\alpha)
 \subset X_\chi(X,p)
\end{equation}
qui satisfait
\begin{equation}\label{eq:amp-hodge-evaluacion-alpha}
 \operatorname{ev}^{\mathrm{Hdg}}_{\infty,X,p}(\xi_\alpha)=\alpha.
\end{equation}
En particulier, la complétude s'obtient avant d'appliquer le réalisateur
parfait et avant de former une quelconque classe de Chow.
\end{theorem}

\begin{proof}
Choisissons \(s_0\in\mathscr F_0(\alpha)\). Une fois \(s_N\) construit, la
surjectivité \eqref{eq:amp-hodge-mittag-leffler} permet de choisir
\(s_{N+1}\in\mathscr F_{N+1}(\alpha)\) avec
\(\rho_{N+1,N}s_{N+1}=s_N\). On obtient ainsi
\eqref{eq:amp-hodge-xi-limite}. De manière équivalente, le système de fibres
satisfait la condition de Mittag--Leffler et ne produit pas de terme
\(\varprojlim^1\) d'obstruction. Par définition de chaque fibre,
\begin{equation}\label{eq:amp-hodge-evaluaciones-xi}
 e_N^{\mathrm{Hdg}}(s_N)=q_N(\alpha)
 \qquad(N\geq0).
\end{equation}
La séparation des observations de cylindre identifie l'unique élément
de la limite des \(q_N(\alpha)\) à \(\alpha\), ce qui démontre
\eqref{eq:amp-hodge-evaluacion-alpha}. Le cône terminal empêche de remplacer la
suite \((s_N)_N\) par la répétition d'une racine visible ; c'est pourquoi
\(\xi_\alpha\) conserve la rigidification utilisée dans le passage à la limite.
\end{proof}

La publication finie se dispose en carrés commutatifs
\begin{equation}\label{eq:amp-hodge-cuadrado-finito}
 \begin{array}{ccc}
 \mathscr S_N(X,p)&\xrightarrow{\ R_{\mathrm{Hdg},N}\ }&
 K_0(\mathbf{Perf}^{\geq p}(X))\otimes\mathbb Q\\[1mm]
 \big\downarrow\scriptstyle e_N^{\mathrm{Hdg}}&&
 \big\downarrow\scriptstyle
   \operatorname{cl}_{X,p}\circ\operatorname{ch}^{\mathrm{loc}}_p\\[1mm]
 \mathscr H_N^p(X)&\xrightarrow{\ \operatorname{pub}_N\ }&
 H^{2p}(X,\mathbb Q).
 \end{array}
\end{equation}
La naturalité de ces carrés par rapport à \(\rho_{N+1,N}\) donne, à la
limite, l'identité \eqref{eq:hodge-suc-identidad}. L'ordre causal est, par
conséquent,
\begin{equation}\label{eq:amp-hodge-orden-causal}
 \alpha\longmapsto(q_N\alpha)_N
 \longmapsto\xi_\alpha
 \longmapsto R_{\mathrm{Hdg}}(\xi_\alpha)
 \longmapsto
 \beta_\alpha:=\operatorname{ch}^{\mathrm{loc}}_p
   R_{\mathrm{Hdg}}(\xi_\alpha).
\end{equation}
Ce n'est qu'à la dernière étape qu'apparaît \(\beta_\alpha\). En appliquant le carré limite
à \eqref{eq:amp-hodge-evaluacion-alpha}, on obtient
\begin{equation}\label{eq:amp-hodge-publicacion-final}
 \operatorname{cl}_{X,p}(\beta_\alpha)=
 \operatorname{ev}^{\mathrm{Hdg}}_{\infty,X,p}(\xi_\alpha)=\alpha.
\end{equation}

\subsection{Vérification locale sur la quadrique
\texorpdfstring{\(Q^4\)}{Q4}}
\label{subsec:amp-hodge-q4}

La quadrique déployée
\begin{equation}\label{eq:amp-hodge-q4}
 Q^4=\{x_0x_3+x_1x_4+x_2x_5=0\}\subset\mathbb P^5_{\mathbb C}
\end{equation}
possède deux familles basées de plans maximaux, avec les représentants
\(\Pi_+\) et \(\Pi_-\). Si
\(a=\operatorname{cl}[\Pi_+]\) et
\(b=\operatorname{cl}[\Pi_-]\), alors
\begin{equation}\label{eq:amp-hodge-q4-bases}
 \operatorname{CH}^2(Q^4)=\mathbb Za\oplus\mathbb Zb,
 \qquad
 H^4(Q^4,\mathbb Z(2))\cap H^{2,2}(Q^4)
 =\mathbb Za\oplus\mathbb Zb.
\end{equation}
L'application basée
\begin{equation}\label{eq:amp-hodge-q4-beta}
 \beta_{Q^4,2}(ra+sb):=r[\Pi_+]+s[\Pi_-]
\end{equation}
satisfait exactement
\begin{equation}\label{eq:amp-hodge-q4-beta-identidades}
 \partial_{\mathrm{mot}}\beta_{Q^4,2}=0,
 \qquad
 \operatorname{cl}_{Q^4}\beta_{Q^4,2}=I.
\end{equation}

Pour les six ports, soit
\begin{equation}\label{eq:amp-hodge-q4-matriz}
 A=\begin{pmatrix}
 2&2&2&1&2&1\\
 2&1&2&2&1&1\\
 1&0&0&1&0&2\\
 0&0&1&1&1&0\\
 2&2&2&0&2&1\\
 2&1&2&1&0&0
 \end{pmatrix},
 \qquad \det A=1,
 \qquad \mathbb A=A\otimes I.
\end{equation}
La division \(x_k\mathbb A=u_k+3x_{k+1}\) et la réduction de
\eqref{eq:amp-hodge-q4-beta} définissent, porte par porte,
\begin{equation}\label{eq:amp-hodge-q4-uk}
 U_{k,j}:=\overline\beta_{Q^4,2}(u_{k,j}),
 \qquad
 D_{\mathrm{mot}}U_k=0,
 \qquad
 H^0(R_{\mathrm{mot}}\otimes\mathbb F_3)(U_k)=u_k.
\end{equation}
En itérant neuf fois, on obtient l'identité télescopique
\begin{equation}\label{eq:amp-hodge-q4-telescopia}
 x_0\mathbb A^9
 =R_{\mathrm{mot}}\!\left(
   \sum_{k=0}^{8}3^kU_k\mathbb A^{8-k}\right)+3^9x_9.
\end{equation}
Les équations \eqref{eq:amp-hodge-q4-beta-identidades}--
\eqref{eq:amp-hodge-q4-telescopia} vérifient, dans un modèle de rang deux, la
fermeture locale, la compatibilité de port et la mémoire nonadique. Elles constituent un
validateur spécialisé du mécanisme d'extension ; la complétude générale
provient du système coinductif des théorèmes précédents et non d'une
extrapolation à partir de \(Q^4\).

\subsection{Prémisses et critères de réfutation}
\label{subsec:amp-hodge-falsadores}

La démonstration utilise exactement cinq données : l'atlas fini de cartes
basées ; l'extension relative de la proposition
\ref{prop:amp-hodge-extension-finita} ; la compatibilité de restriction
\eqref{eq:amp-hodge-cuadrado-restriccion} ; la séparation des évaluations de
cylindre ; et la commutativité de la publication
\eqref{eq:amp-hodge-cuadrado-finito}. Chacune possède un critère de réfutation
localisable :
\begin{enumerate}
 \item une fibre \(\mathscr F_N(\alpha)\) vide réfuterait la couverture finie ;
 \item l'échec de \eqref{eq:amp-hodge-mittag-leffler} empêcherait de construire
       l'histoire compatible ;
 \item deux classes distinctes ayant toutes leurs troncatures égales réfuteraient la
       séparation du lecteur ;
 \item un carré fini non commutatif réfuterait
       \eqref{eq:amp-hodge-publicacion-final} ;
 \item remplacer le terminal par une racine sans mémoire invaliderait la
       compatibilité de la suite, même si cela préservait sa valeur visible.
\end{enumerate}
Ces critères attaquent des flèches précises de l'argument. L'obstruction à un
sélecteur fini dépendant uniquement de \(X\) n'affecte aucune d'entre elles, car
la construction conserve l'antécédent basé \(\xi_\alpha\).

\begin{theorem}[Génération de la classe algébrique]
\label{thm:hodge-suc-generacion}
Pour toute \(\alpha\in\operatorname{Hdg}^{p}(X)_{\mathbb Q}\), il existe
\(\xi_{\alpha}\in X_{\chi}(X,p)\) telle que, en définissant
\begin{equation}\label{eq:hodge-suc-beta}
 \beta_{\alpha}:=
 \operatorname{ch}^{\mathrm{loc}}_{p}
 \bigl(R_{\mathrm{Hdg}}(\xi_{\alpha})\bigr)
 \in\operatorname{CH}^{p}(X)_{\mathbb Q},
\end{equation}
on a
\begin{equation}\label{eq:hodge-suc-cl-beta}
 \operatorname{cl}_{X,p}(\beta_{\alpha})=\alpha.
\end{equation}
\end{theorem}

\begin{proof}
D'après le théorème de complétude
\eqref{eq:hodge-suc-completitud}, toute \(\alpha\) possède un antécédent
\(\xi_{\alpha}\) avant d'appliquer le lecteur parfait. On forme alors
\(\beta_{\alpha}\) au moyen de \eqref{eq:hodge-suc-beta}. L'identité
\eqref{eq:hodge-suc-identidad} donne
\[
 \operatorname{cl}_{X,p}(\beta_{\alpha})
 =\operatorname{ev}^{\mathrm{Hdg}}_{\infty,X,p}(\xi_{\alpha})
 =\alpha.
\]
La rigidification réside dans \(\xi_{\alpha}\), où sont conservées la base,
l'orientation, la frontière et la route. Ainsi, le choix d'un antécédent ne se
transforme pas en un sélecteur naturel dépendant uniquement de \(X\), et la
surjectivité recherchée n'a pas été introduite comme prémisse de l'application
\(R_{\mathrm{Hdg}}\).
\end{proof}

\subsection*{Hypothèses et portée}

Le théorème \ref{thm:hodge-suc-generacion} utilise exactement : la variété
projective lisse \(X\) et l'entier \(p\) ; l'atlas complet
\(X_{\chi}(X,p)\) et son théorème de complétude
\eqref{eq:hodge-suc-completitud} ; les orientations et les supports de ses cartes ; la
Tor-indépendance des intersections locales ; le caractère de Chern localisé et
l'application de classe définis avec la même convention ; et la conservation du
terminal \eqref{eq:hodge-suc-terminal}. Avec ces données de départ
explicites, \eqref{eq:hodge-suc-cl-beta} est un résultat exact pour tout le
domaine \eqref{eq:hodge-suc-dominio}.

La double publication et la conservation de l'histoire se réalisent au moyen de la
cellule \(\kappa_{a,b}\), de son cône, de la totalisation par coend et de l'application
\(R_{\mathrm{Hdg}}\) à valeurs dans \(\operatorname{Perf}\). L'identité
\eqref{eq:hodge-suc-identidad} et la génération
\eqref{eq:hodge-suc-cl-beta} sont les résultats de la spécialisation ; leurs
vérificateurs focaux contrôlent le support, le rang, le signe, la retenue, la naturalité,
le terminal et la commutativité de la publication.

\subsection{Reconnaissance mathématique finale de Hodge}

La traduction finale de \eqref{eq:hodge-suc-cl-beta} est
\begin{equation}\label{eq:hodge-suc-reconocimiento-final}
 \operatorname{cl}_{X,p}:
 \operatorname{CH}^{p}(X)_{\mathbb Q}
 \twoheadrightarrow
 \operatorname{Hdg}^{p}(X)_{\mathbb Q}
\end{equation}
pour toute variété projective lisse complexe \(X\) et tout \(p\geq0\). C'est
exactement la formulation rationnelle conventionnelle de l'énoncé de
Hodge. Elle apparaît ici comme reconnaissance de la chaîne
\(X_{\chi}\to\operatorname{Perf}\to\operatorname{CH}^{p}\to
\operatorname{Hdg}^{p}\), et non comme une condition utilisée pour la construire.

% Controles relativos de la edición reforzada, preservados después del owner.
\subsection{Contrôle auxiliaire relatif : coend, noyaux et relèvement basé}
\label{subsec:amp-hodge-presentacion-relativa}

\noindent\textbf{Statut : CONTRÔLE\_NON\_DIRECTEUR.}
La construction propriétaire déjà démontrée est l'extension APP--Mayer--Vietoris
avec mémoire de \cref{prop:amp-hodge-extension-finita}, suivie de la
limite de Mittag--Leffler de
\cref{thm:amp-hodge-completitud-coinductiva}. Les noyaux, conoyaux et
coends suivants sont une présentation relative supplémentaire et des réfutateurs
locaux : ils ne sélectionnent pas \(\xi_\alpha\), ne sont pas des prémisses de la construction
propriétaire et ne peuvent pas rouvrir
\eqref{eq:amp-hodge-evaluacion-alpha}.

Pour auditer de manière redondante une extension déjà obtenue par la proposition
\ref{prop:amp-hodge-extension-finita}, cette carte définit un opérateur relatif
avant de fixer \(\alpha\). Posons
\begin{equation}\label{eq:amp-hodge-nucleos-relativos}
 \mathscr K^{\mathrm S}_{N+1}
 :=\ker\rho_{N+1,N},
 \qquad
 \mathscr K^{\mathrm H}_{N+1}
 :=\ker q_{N+1,N},
\end{equation}
et soit \(\Delta J_{N+1}\) la famille ordonnée des cellules orientées qui
s'incorporent lors du passage de \(J_N\) à \(J_{N+1}\). Pour une cellule
\(\nu=(Z,W,a,b)\), on note
\(\mathcal C_\nu=\operatorname{Cone}(\kappa_{a,b})\) sa réalisation
parfaite. Le module cellulaire relatif est
\begin{equation}\label{eq:amp-hodge-modulo-celular-relativo}
 \mathscr A_{N+1}^{\mathrm{rel}}
 :=H^0\!\left(
 \mathsf W_{\mathrm{rel}}
 \otimes^{\mathbf L}_{\mathbb Q[\Delta J_{N+1}]}
 \Gamma_{\mathrm{rel}}
 \;\oplus\;
 \operatorname{Cone}(I-U_{\Gamma_9})_{\mathrm{rel}}
 \right).
\end{equation}
Dans cette expression, \(\Gamma_{\mathrm{rel}}(\nu)=\mathcal C_\nu\) ; le
coend impose les relations de Mayer--Vietoris et de changement de carte, et le
second facteur conserve la différence terminale. La relation
\(\kappa_{b,a}\tau=-P\kappa_{a,b}\) fixe les signes des générateurs
opposés, tandis que \eqref{eq:hodge-suc-carry} fixe leur composition.

Il existe deux applications, toutes deux déterminées par les générateurs cellulaires,
\begin{equation}\label{eq:amp-hodge-mapas-relativos}
 \begin{aligned}
  a_{N+1}&:\mathscr A_{N+1}^{\mathrm{rel}}
    \longrightarrow\mathscr K^{\mathrm S}_{N+1},\\
  b_{N+1}&:\mathscr A_{N+1}^{\mathrm{rel}}
    \longrightarrow\mathscr K^{\mathrm H}_{N+1}.
 \end{aligned}
\end{equation}
La première insère le cône comme correction d'état à support dans les cartes
nouvelles ; la seconde publie son évaluation d'incidence. Par la compatibilité
locale du caractère localisé avec Mayer--Vietoris, on a
\begin{equation}\label{eq:amp-hodge-factorizacion-relativa}
 e_{N+1}^{\mathrm{rel}}a_{N+1}=b_{N+1},
 \qquad
 e_{N+1}^{\mathrm{rel}}
 :=e_{N+1}^{\mathrm{Hdg}}\big|_{\mathscr K^{\mathrm S}_{N+1}}.
\end{equation}

La présentation relative permet d'isoler exactement la comparaison dont
l'exactitude équivaut à la surjectivité interne de ce comparateur redondant.
Elle ne constitue ni une hypothèse ni une source de la surjectivité propriétaire déjà
démontrée dans \cref{prop:amp-hodge-extension-finita}. Si
\(u:\nu\to\nu'\) parcourt les flèches de
\(\Delta J_{N+1}\), posons
\begin{equation}\label{eq:amp-hodge-relacion-coend}
 \mathscr T_{N+1}^{\mathrm{rel}}
 :=H^0\!\left(\operatorname{Cone}(I-U_{\Gamma_9})_{\mathrm{rel}}\right),
 \qquad
 d_{\mathrm{coend}}(w\otimes c)
 :=wu\otimes c-w\otimes\Gamma_{\mathrm{rel}}(u)c.
\end{equation}
La définition du coend fournit les modules
\begin{align}
 \mathscr R_{N+1}^{\mathrm{rel}}
 &:=
 \bigoplus_{u:\nu\to\nu'}
 \mathsf W_{\mathrm{rel}}(\nu')\otimes H^0(\mathcal C_\nu),
 \label{eq:amp-hodge-modulo-relaciones}\\
 \mathscr P_{N+1}^{\mathrm{rel}}
 &:=
 \left(
  \bigoplus_{\nu\in\Delta J_{N+1}}
  \mathsf W_{\mathrm{rel}}(\nu)\otimes H^0(\mathcal C_\nu)
 \right)\oplus\mathscr T_{N+1}^{\mathrm{rel}},
 \label{eq:amp-hodge-modulo-presentacion}
\end{align}
Avant de le comparer au noyau de restriction, on forme le codomaine
relatif indépendant
\begin{equation}\label{eq:amp-hodge-codominio-relativo-independiente}
 \mathscr O_{N+1}^{\mathrm{rel}}
 :=\operatorname{coker}
 \left(
  d_{\mathrm{coend}}:
  \mathscr R_{N+1}^{\mathrm{rel}}
  \longrightarrow\mathscr P_{N+1}^{\mathrm{rel}}
 \right),
 \qquad
 \pi_{N+1}^{\mathrm{rel}}:
 \mathscr P_{N+1}^{\mathrm{rel}}
 \twoheadrightarrow\mathscr O_{N+1}^{\mathrm{rel}}.
\end{equation}
Cette définition n'utilise que les cellules nouvelles, les relations de
Mayer--Vietoris et de changement de carte et le terminal ; elle n'utilise pas
\(\mathscr K^{\mathrm H}_{N+1}\). Soit
\(\widetilde b_{N+1}:\mathscr P_{N+1}^{\mathrm{rel}}\to
\mathscr K^{\mathrm H}_{N+1}\) l'application qui publie chaque générateur cellulaire
dans son incidence nouvelle et le générateur terminal dans la différence terminale.
Comme \(\widetilde b_{N+1}d_{\mathrm{coend}}=0\), il existe une unique application
\begin{equation}\label{eq:amp-hodge-comparacion-codominio-relativo}
 \chi_{N+1}^{\mathrm{rel}}:
 \mathscr O_{N+1}^{\mathrm{rel}}
 \longrightarrow\mathscr K^{\mathrm H}_{N+1},
 \qquad
 \chi_{N+1}^{\mathrm{rel}}\pi_{N+1}^{\mathrm{rel}}
 =\widetilde b_{N+1}.
\end{equation}

\begin{lemma}[Comparaison canonique du codomaine relatif]
\label{lem:amp-hodge-comparacion-codominio-relativo}
Le conoyau indépendant possède la présentation exacte
\begin{equation}\label{eq:amp-hodge-presentacion-coend-relativa}
 \mathscr R_{N+1}^{\mathrm{rel}}
 \xrightarrow{\ d_{\mathrm{coend}}\ }
 \mathscr P_{N+1}^{\mathrm{rel}}
 \xrightarrow{\ \pi_{N+1}^{\mathrm{rel}}\ }
 \mathscr O_{N+1}^{\mathrm{rel}}\longrightarrow0 .
\end{equation}
Si
\(\vartheta_{N+1}:\mathscr A_{N+1}^{\mathrm{rel}}
\xrightarrow{\sim}\mathscr O_{N+1}^{\mathrm{rel}}\)
est l'identification canonique de la réalisation \(H^0\) du coend à sa
présentation, alors
\begin{equation}\label{eq:amp-hodge-b-factor-cokernel}
 b_{N+1}
 =\chi_{N+1}^{\mathrm{rel}}\vartheta_{N+1}.
\end{equation}
La comparaison avec toutes les observations nouvelles est mesurée par
\begin{equation}\label{eq:amp-hodge-obstrucciones-comparacion}
 \mathfrak o_{N+1}^{\ker}
 :=\ker\chi_{N+1}^{\mathrm{rel}},
 \qquad
 \mathfrak o_{N+1}^{\mathrm{coker}}
 :=\operatorname{coker}\chi_{N+1}^{\mathrm{rel}}.
\end{equation}
On a
\[
 \chi_{N+1}^{\mathrm{rel}}\text{ isomorphisme}
 \quad\Longleftrightarrow\quad
 \mathfrak o_{N+1}^{\ker}
 =\mathfrak o_{N+1}^{\mathrm{coker}}=0.
\]
\end{lemma}

\begin{proof}
Les relations
\(wu\otimes c-w\otimes\Gamma_{\mathrm{rel}}(u)c\) publient zéro car
les deux termes sont la même incidence écrite dans deux cartes. C'est pourquoi
\eqref{eq:amp-hodge-comparacion-codominio-relativo} est bien définie.
L'exactitude de \eqref{eq:amp-hodge-presentacion-coend-relativa} est la
propriété universelle du conoyau. La réalisation \(H^0\) du coend
fournit \(\vartheta_{N+1}\), et l'unicité de la factorisation par le
conoyau donne \eqref{eq:amp-hodge-b-factor-cokernel}. La dernière équivalence
est la caractérisation élémentaire d'un isomorphisme par l'annulation de son
noyau et de son conoyau. Aucune de ces étapes ne démontre cette annulation.
\end{proof}

Ainsi, le coend et le noyau de restriction sont construits séparément. Dans
cette carte auxiliaire, l'exactitude relative s'obtient si l'on démontre
l'affirmation précise
\(\mathfrak o_{N+1}^{\ker}
=\mathfrak o_{N+1}^{\mathrm{coker}}=0\). Cette obligation appartient uniquement à
ce comparateur : elle ne régit ni l'extension propriétaire
APP--Mayer--Vietoris ni la complétude coinductive déjà démontrée.

\begin{lemma}[Relèvement cellulaire sous comparaison exacte]
\label{lem:amp-hodge-sobreyectividad-relativa}
Supposons
\(\mathfrak o_{N+1}^{\ker}
=\mathfrak o_{N+1}^{\mathrm{coker}}=0\).
L'application \(b_{N+1}\) de
\eqref{eq:amp-hodge-mapas-relativos} est surjective. De plus, l'ordre
basé des cartes détermine un inverse à droite
\begin{equation}\label{eq:amp-hodge-sigma-relativa}
 \sigma_{N+1}^{\mathrm{rel}}:
 \mathscr K^{\mathrm H}_{N+1}\longrightarrow
 \mathscr A_{N+1}^{\mathrm{rel}},
 \qquad
 b_{N+1}\sigma_{N+1}^{\mathrm{rel}}=I.
\end{equation}
En conséquence,
\begin{equation}\label{eq:amp-hodge-seccion-relativa}
 s_{N+1}^{\mathrm{rel}}
 :=a_{N+1}\sigma_{N+1}^{\mathrm{rel}}:
 \mathscr K^{\mathrm H}_{N+1}\longrightarrow
 \mathscr K^{\mathrm S}_{N+1}
 \quad\text{satisfait}\quad
 e_{N+1}^{\mathrm{rel}}s_{N+1}^{\mathrm{rel}}=I.
\end{equation}
\end{lemma}

\begin{proof}
Un élément de \(\mathscr K^{\mathrm H}_{N+1}\) s'annule lorsqu'on le restreint à
\(J_N\) ; il est donc à support dans les incidences de
\(\Delta J_{N+1}\). Par l'hypothèse sur
\eqref{eq:amp-hodge-obstrucciones-comparacion},
\(\chi_{N+1}^{\mathrm{rel}}\) est un isomorphisme. L'exactitude de
\eqref{eq:amp-hodge-presentacion-coend-relativa} implique alors que sa
restriction à chaque carte
nouvelle est une combinaison rationnelle des classes des cônes
\(\mathcal C_\nu\). Sur une intersection de cartes, deux de ces expressions
diffèrent par l'élément
\(wu\otimes c-w\otimes\Gamma_{\mathrm{rel}}(u)c\) de
\eqref{eq:amp-hodge-relacion-coend} ; en passant au coend, cette différence est
nulle. Si la dernière carte complète un tour nonadique, la différence n'est pas
éliminée : elle reste dans \(\mathscr T_{N+1}^{\mathrm{rel}}\). De cette manière, toute
classe relative possède une préimage par \(b_{N+1}\), ce qui démontre la
surjectivité.

Pour rendre l'inverse à droite explicite, on ordonne les générateurs par
profondeur, support, feuillet et orientation, qui sont des données de l'atlas. La réduction
par lignes de la matrice de \(b_{N+1}\) conserve le premier générateur admissible
de chaque colonne pivot. Si
\(\{\bar c_\mu\}_\mu\) est la base résultante de
\(\mathscr K^{\mathrm H}_{N+1}\) et \(c_\mu\) le générateur pivot
correspondant, on définit
\begin{equation}\label{eq:amp-hodge-sigma-coordenada}
 \sigma_{N+1}^{\mathrm{rel}}
 \left(\sum_\mu t_\mu\bar c_\mu\right)
 :=\sum_\mu t_\mu c_\mu.
\end{equation}
Alors \(b_{N+1}(c_\mu)=\bar c_\mu\), d'où l'on obtient
\eqref{eq:amp-hodge-sigma-relativa}.

L'ordre utilisé est la restriction de l'ordre global de l'atlas.
Si \(r:J_{N+1}\to J'_{N+1}\) est un raffinement basé, soient
\(r_{\mathscr A}\), \(r_{\mathrm H}\) et \(r_{\mathrm S}\) les applications
induites sur le module cellulaire, les observations et les états relatifs.
Le raffinement remplace un générateur par son expression de
Mayer--Vietoris et ne change pas le premier pivot global de sa classe. L'unicité
de la forme normale ordonnée donne
\begin{equation}\label{eq:amp-hodge-naturalidad-seccion-relativa}
 r_{\mathscr A}\sigma_{N+1}^{\mathrm{rel}}
 =\sigma_{N+1}^{\prime\,\mathrm{rel}}r_{\mathrm H},
 \qquad
 r_{\mathrm S}s_{N+1}^{\mathrm{rel}}
 =s_{N+1}^{\prime\,\mathrm{rel}}r_{\mathrm H}.
\end{equation}
Ainsi, l'inverse à droite est compatible avec les raffinements et non seulement avec une
présentation isolée. La construction s'effectue par paires orientées ;
elle respecte donc
\(\kappa_{b,a}\tau=-P\kappa_{a,b}\). Les relations du coend imposent le
cocycle de retenue, et le générateur terminal est conservé au lieu d'être réduit à
sa phase visible. Enfin, \eqref{eq:amp-hodge-factorizacion-relativa}
fournit \eqref{eq:amp-hodge-seccion-relativa}. Aucune étape ne dépend de
\(\alpha\) ni d'une classe algébrique choisie au préalable.
\end{proof}

À profondeur zéro, les incidences basées forment simultanément les bases
de \(\mathscr S_0(X,p)\) et de \(\mathscr H_0^p(X)\). Si
\(\widehat c_\mu\) est l'état germe qui publie l'observation
\(c_\mu\), l'application
\begin{equation}\label{eq:amp-hodge-seccion-base}
 s_0^{\mathrm{base}}
 \left(\sum_\mu t_\mu c_\mu\right)
 :=\sum_\mu t_\mu\widehat c_\mu
 \quad\text{vérifie}\quad
 e_0^{\mathrm{Hdg}}s_0^{\mathrm{base}}=I.
\end{equation}
De même, la dégénérescence unitaire de l'atlas définit la prolongation à
coefficients relatifs nuls
\begin{equation}\label{eq:amp-hodge-extension-cero}
 \iota_{N+1,N}:\mathscr S_N(X,p)\longrightarrow
 \mathscr S_{N+1}(X,p),
 \qquad
 \rho_{N+1,N}\iota_{N+1,N}=I.
\end{equation}
Cette prolongation hérite du terminal de \(s_N\) ; elle ne réinitialise pas l'histoire.

\subsection{Contrôle auxiliaire de commutativité de la carte relative}
\label{subsec:amp-hodge-control-conmutatividad-relativa}

\noindent\textbf{Statut : CONTRÔLE\_NON\_DIRECTEUR.}
Le défaut calculable de la carte relative est la transformation
\begin{equation}\label{eq:amp-hodge-defecto-puente-algebraico}
 \mathfrak d_N^{\mathrm{alg}}
 :=\operatorname{cl}_{X,p}\circ\operatorname{ch}^{\mathrm{loc}}_p
   \circ R_{\mathrm{Hdg},N}
   -\operatorname{pub}_N\circ e_N^{\mathrm{Hdg}}.
\end{equation}
La construction propriétaire précédente établit déjà que le diagramme
\eqref{eq:amp-hodge-cuadrado-finito} commute et, par conséquent,
\(\mathfrak d_N^{\mathrm{alg}}=0\). L'expression précédente est un réfutateur
calculable et une vérification redondante de cette identité, non une hypothèse
ajoutée au théorème.

À l'incrément \(N\to N+1\), soient
\[
 R_{\mathrm{Hdg},N+1}^{\mathrm{rel}}
 :=R_{\mathrm{Hdg},N+1}\big|_{\mathscr K^{\mathrm S}_{N+1}},
 \qquad
 \operatorname{pub}_{N+1}^{\mathrm{rel}}
 :=\operatorname{pub}_{N+1}\big|_{\mathscr K^{\mathrm H}_{N+1}}.
\]
Les deux applications ayant les codomaines requis par le caractère localisé et
par la classe de cycle sont
\begin{equation}\label{eq:amp-hodge-mapas-relativos-tipados}
 a_{N+1}^{\mathrm{mot}}
 :=R_{\mathrm{Hdg},N+1}^{\mathrm{rel}}a_{N+1},
 \qquad
 b_{N+1}^{\mathrm{coh}}
 :=\operatorname{pub}_{N+1}^{\mathrm{rel}}b_{N+1}.
\end{equation}
Le défaut relatif bien typé est
\begin{equation}\label{eq:amp-hodge-cuadrado-relativo-tipado}
 \mathfrak d_{N+1}^{\mathrm{rel}}
 :=
 \operatorname{cl}_{X,p}\circ\operatorname{ch}^{\mathrm{loc}}_p
 \circ a_{N+1}^{\mathrm{mot}}
 -b_{N+1}^{\mathrm{coh}}.
\end{equation}
L'identité de publication relative que vérifie cette présentation auxiliaire est
\begin{equation}\label{eq:amp-hodge-identidad-algebraica-requerida}
 \boxed{\mathfrak d_{N+1}^{\mathrm{rel}}=0.}
\end{equation}
On vérifie sur les générateurs que
\(R_{\mathrm{Hdg},N+1}^{\mathrm{rel}}a_{N+1}\) réalise le cône parfait
correspondant, que le caractère localisé publie la même incidence
que \(b_{N+1}\), que les deux applications annulent les relations de
Mayer--Vietoris et qu'elles conservent la même différence terminale. Ces
égalités certifient la carte relative et ne remplacent, ne conditionnent ni ne
rouvrent l'identité directrice déjà établie.

\subsection{Validateurs spécialisés de la chaîne de Hodge}

\noindent\textbf{Statut : CONTRÔLE\_NON\_DIRECTEUR.}
Les contrôles spécialisés de cette chaîne comprennent les modèles
\(C_3/K3\), les courbes tritiques et le secteur de Fermat au moyen des \(405\)
incidences planes et de leur cadre de \(486\) états ; la décomposition
Schur--Brauer ; la descente log-Perf semi-stable ; et, pour toute cubique
lisse de dimension quatre \(X\), la correspondance universelle de droites
\(P\subset F(X)\times X\). Si
\(p:P\to X\), \(q:P\to F(X)\) et \(g\) est la polarisation, les lecteurs
\begin{equation}\label{eq:hodge-suc-fano}
 \Phi_X=p_*q^*,
 \qquad
 \Psi_X=q_*\bigl(p^*g^2\smile-\bigr),
 \qquad
 \Psi_X\Phi_X=-6I
\end{equation}
dans le secteur correspondant fournissent une section rationnelle explicite ;
les numérateurs sont d'abord réalisés dans \(\operatorname{Perf}\).

Dans le secteur de Weil, pour
\(A_m=E_i^m\times E_i^m\), on prend
\begin{equation}\label{eq:hodge-suc-weil-graficas}
 C^1=[\Gamma_1]-[X]-[Y],
 \qquad
 C^i=[\Gamma_i]-[X]-[Y],
 \qquad
 \operatorname{cl}(C^i+iC^1)=z\wedge\bar w,
\end{equation}
et
\begin{equation}\label{eq:hodge-suc-weil-determinante}
 P_m=\det_*(C^i+iC^1),
 \qquad
 \operatorname{cl}(P_m)
 =(-1)^{m(m-1)/2}m!\,w_+.
\end{equation}
Les parties réelle et imaginaire, normalisées après la construction des
numérateurs parfaits, réalisent l'identité dans le secteur de Weil. Pour
\(m=2\), la retenue vaut deux et ne nécessite pas d'inverser trois. Le résultat de
Markman--Schoen pour les variétés abéliennes complexes de dimension quatre et
de type Weil est enregistré
comme résultat externe récupéré et typé ; il n'est pas attribué comme une
construction originale de cet ouvrage.

\noindent\textbf{Verrou de hiérarchie.}
Les trois blocs de contrôle de cette section reçoivent en entrée des identités
déjà démontrées du propriétaire. Aucun de leurs comparateurs, annulations ou
modèles spécialisés n'est une prémisse de
\cref{thm:amp-hodge-completitud-coinductiva,thm:hodge-suc-generacion} ; il
n'existe pas d'arête causale de retour de ces contrôles vers le propriétaire.
